\section{Limitations}
\label{sec:limitations}

Stating limitations this precisely is itself part of the
contribution argued in \S\ref{sec:related-work}: a reputation
aggregator's side-by-side verdicts carry no equivalent account of where
they might be wrong, which makes them fast to consult but hard to
audit. Every limitation below is instead diagnosed to a specific,
verifiable cause rather than stated generically, so that each one is
either already mitigated concretely or scoped as a precise next step
(\S\ref{sec:future-work}) rather than an open-ended caveat.

\begin{enumerate}

\item \label{lim:n3} \textbf{Small validation set (N=3).} Three families is too few
  for a statistically powered claim, and is the point most deserving of
  scrutiny in this work. Three concrete mitigations: ground truth is
  function-level, Ghidra-verified manual reverse engineering per sample,
  a materially higher bar than a larger but shallower set would carry;
  a full audit of \code{static/scripts/} specifically for
  sample-specific hardcoding disguised as generic logic found and
  removed four instances (a literal installer-helper filename, a
  hardcoded encrypted-payload extension, product strings baked into
  filesystem-detection logic, and a ``generic'' decoder that was
  actually one sample's exact XOR key), each replaced with a
  class-level equivalent and re-verified not to regress any validated
  sample; and RoningLoader's resubmission loop (\S\ref{sec:resubmission})
  independently scores 26 additional components against the same
  ground truth, real additional evaluation surface rather than a
  statistical trick. Extending to further families remains the
  highest-value follow-up (\S\ref{sec:future-work}).

\item \label{lim:attck-version} \textbf{ATT\&CK version is pinned, and
  CAPE's own signature tags are not.} All technique IDs in this pipeline, ground truth, and
  Navigator layers target MITRE ATT\&CK v19.2. CAPE's community
  signatures are frozen to whichever version they were last written
  against, so signature-sourced tags are remapped to their current
  identity at curation time (\S\ref{sec:dynamic}) rather than trusted
  verbatim. Migrating ground truth from the previously-unpinned v14
  identities to v19 is what surfaced the
  \technique{T1070}/\technique{T1685.005} confidence-calibration bug
  detailed in \S\ref{sec:fp-audit}, which a version left unpinned would
  have masked indefinitely.

\item \label{lim:gt-table} \textbf{Ground-truth extraction depends on
  table completeness.}
  A markdown-table-based extractor is only as complete as the table
  itself, even when a report's prose says more.
  \S\ref{sec:case-wsnake} is a concrete instance, fixed at the source
  document rather than worked around in the extractor.

\item \label{lim:containment} \textbf{Containment-first sandboxing
  limits network/C2 observation.} CAPE detonates against INetSim's simulated network with
  no real egress, by design. Several missed techniques across all three
  samples cluster around exfiltration and proxy/tunneling behavior
  (\technique{T1041}, \technique{T1090}, \technique{T1571/T1572},
  \technique{T1046}) that a simulated network cannot fully elicit. This
  is a structural property of a containment-first research environment,
  not a defect to fix.

\item \label{lim:key-ceiling} \textbf{Static key/config recovery has a
  precise ceiling.}
  Single-byte XOR is exhaustively brute-forceable (128 keys), and was
  used to recover RoningLoader's C2 IP with no prior key knowledge, but
  only after two independent false-positive filters were added
  (\S\ref{sec:case-xor}). By contrast, a separate payload's RC4 key,
  independently verified via manual Ghidra RE, is assembled at
  \emph{runtime} from three register-arithmetic-combined constants and
  never exists as a contiguous byte sequence in the file: no static
  byte-pattern method, however exhaustive, can recover a key that is
  never stored as bytes. The boundary is precisely whether the key
  exists as bytes on disk, not a general statement that ``encryption is
  hard.''

\item \label{lim:tool-constraints} \textbf{Tool-level constraints.} FLOSS's emulation-based string
  decoding timed out on Akira, the most anti-analysis-heavy sample,
  logged explicitly rather than silently absent. AES/ChaCha20 recovery
  only succeeds against a zero nonce/IV. Any new brute-force-based
  recovery capability is now, by explicit policy, validated against a
  second sample before being trusted, after the XOR key false-positive
  cascade (\S\ref{sec:case-xor}) was caught this way rather than on
  first success.

\item \label{lim:reproducibility} \textbf{Reproducibility.} \S\ref{sec:deployment} describes the
  full \code{docker-compose} deployment; the one gap worth restating
  here is that the \code{sandbox} profile's Windows guest VM needs
  one-time manual host provisioning (libvirt/KVM), documented with its
  own troubleshooting notes, since it cannot be created by
  \code{docker-compose} alone.

\item \label{lim:rule-coverage} \textbf{Static rule coverage is
  deliberately narrow.} Roughly 85
  of 474 Windows-relevant MITRE techniques have a dedicated static
  rule, the scoping rationale detailed in \S\ref{sec:coverage}.

\item \label{lim:network-attribution} \textbf{Dynamic network capture
  has no per-process attribution.}
  CAPE's \code{network.hosts} list reflects the whole sandbox VM's
  traffic, not just the sample's own process, so the guest OS's own
  background telemetry can appear indistinguishably from the sample's
  real connections (\S\ref{sec:case-ioc-noise}). The mitigation applied
  is a narrow, individually-confirmed allowlist of the specific IPs
  actually observed contaminating these three samples' detonations, not
  a general ASN or reputation filter; a newly detonated sample could
  surface the same class of contamination under a different IP the
  current allowlist does not yet cover.

\end{enumerate}
